Het Network Time Protocol\index{Network Time Protocol} haalt de datum en tijd op vanaf servers op het Internet die direct of indirect verbonden zijn met een atoomklok. Via deze methode kunnen alle systemen op aarde dezelfde referentie tijd (UTC\index{UTC}, Coordinated Universal Time\index{Coordinated Universal Time}) gebruiken. UTC kent geen zomer- en wintertijd. Er is dus altijd kennis van de tijdzone nodig om UTC om te rekenen naar lokale tijd. Het nul-punt van UTC ligt gelijk aan de GMT.

Het \texttt{systemd} systeem heeft zijn eigen implementatie van een NTP\index{NTP}-server. De daemon heet \texttt{systemd-timesyncd}\index{systemd-timesyncd}. Controlleer of je systeem gebruik maakt van deze NTP-server:
\begin{lstlisting}
systemctl status systemd-timesyncd
\end{lstlisting}

